\section{Methodology}
\label{sec:methodology}

This section details the experimental methodology employed in the PW-2026-005 study to discriminate between fractional, delayed, and latent ecological memory mechanisms in predator-prey systems with strong Allee effects. The experimental design integrates recent advances in fractional-order system identification (2024-2025) and adheres to the dimensional consistency requirements of the Caputo fractional derivative framework.

\subsection{Experimental System}
We consider the dimensionally consistent Caputo predator-prey model with strong Allee effect:
\begin{equation}
\label{eq:system}
\tau_0^{\alpha-1} {}_0D_t^\alpha \mathbf{z}(t) = \mathbf{f}(\mathbf{z}(t);\boldsymbol{\theta}) + \mathbf{B}\mathbf{u}(t),
\end{equation}
where $\mathbf{z}(t) = [x(t), y(t)]^\top$ denotes the prey ($x$) and predator ($y$) populations, $\tau_0$ is a characteristic time scale, $\alpha \in (0,1]$ is the fractional order, ${}_0D_t^\alpha$ denotes the Caputo fractional derivative of order $\alpha$, $\mathbf{f}$ represents the nonlinear dynamics, $\boldsymbol{\theta}$ are model parameters, $\mathbf{B}$ is the input matrix, and $\mathbf{u}(t)$ is the control input. The output equation is $\mathbf{y}(t) = \mathbf{C}\mathbf{z}(t)$.

\subsection{Memory Hypotheses}
Three competing memory hypotheses are formulated and tested:
\begin{enumerate}
    \item \textbf{Fractional Memory}: Memory kernel follows a power-law decay $\frac{t^{\alpha-1}}{\Gamma(\alpha)}$.
    \item \textbf{Delayed Memory}: Discrete time-delay representation with distributed or discrete delays.
    \item \textbf{Latent Memory}: Finite-dimensional latent state approximation (e.g., linear compartmental model).
\end{enumerate}

\subsection{Experimental Protocol}
The experiment comprises three sequential phases:

\subsubsection{Phase 1: Baseline Characterization}
\begin{itemize}
    \item Establish the coexistence equilibrium $(x^*, y^*) = \left(\frac{2}{3}, \frac{2(2-3A)}{9A}\right)$ for Allee threshold $A=0.3$.
    \item Verify local asymptotic stability for $\alpha < \alpha^*(A) \approx 0.969$ via linearized fractional-order dynamics.
    \item Apply infinitesimal perturbations to estimate natural decay rates and baseline noise floor.
\end{itemize}

\subsubsection{Phase 2: Perturbation Experiments}
We apply three classes of controlled perturbations, each designed to excite distinct memory signatures:
\begin{enumerate}
    \item \textbf{Prey Pulses}: Brief increases in prey population ($\Delta x = 0.2$) via transient feeding.
    \item \textbf{Predator Pulses}: Brief increases in predator population ($\Delta y = 0.15$) via transient predator introduction.
    \item \textbf{Chirp Signals}: Frequency-swept inputs $u(t) = A \sin(\phi(t))$ with instantaneous frequency $\phi'(t)$ logarithmically increasing from $10^{-3}$ to $10^{1}$ rad/$t_0$ over $T=30\tau_0$, to estimate frequency response.
\end{enumerate}
All perturbations are designed to respect safety constraints: $0 < x,y < 10$, $\|\mathbf{u}(t)\| \leq 1.0$, and $x(t) > A$.

\subsubsection{Phase 3: Measurement Protocol}
\begin{itemize}
    \item High-frequency sampling immediately post-perturbation ($\Delta t = 0.01\tau_0$) for $t \in [0, 5\tau_0]$.
    \item Extended monitoring to capture long-term memory effects up to $T = 50\tau_0$ with $\Delta t = 0.1\tau_0$.
    \item Dual-species monitoring via non-invasive imaging to capture cross-species memory effects.
    \item Each experimental condition is replicated $n=5$ times to ensure statistical significance.
\end{itemize}

\subsection{Fractional Memory Signatures}
Discrimination relies on temporal and frequency-domain signatures:

\subsubsection{Temporal Signature}
The system's impulse response decay distinguishes memory types:
\begin{itemize}
    \item \textbf{Fractional}: Power-law decay $y(t) \sim t^{-\alpha}$.
    \item \textbf{Delayed}: Exponential decay with time delay $y(t) \sim e^{-(t-\tau)/\tau_d}$.
    \item \textbf{Latent}: Multi-exponential decay $y(t) \sim \sum_{i=1}^n a_i e^{-t/\tau_i}$.
\end{itemize}

\subsubsection{Frequency Domain Signature}
The frequency response $G_\alpha(i\omega)$ yields:
\begin{itemize}
    \item \textbf{Fractional}: Constant phase angle $-(1-\alpha)\pi/2$ in intermediate frequencies (Bode phase plateau).
    \item \textbf{Delayed}: Linear phase shift $-\omega\tau$.
    \item \textbf{Latent}: Rational function approximation with bounded phase variation.
\end{itemize}

\subsection{Incorporation of Recent Findings (2024-2025)}
Recent advances in fractional-order ecological modeling are integrated:
\begin{itemize}
    \item Mittag-Leffler kernel fractional derivatives provide superior fit to ecological memory data versus power-law kernels (see \citealp{mittag2024ecological}).
    \item Distributed-order fractional derivatives capture heterogeneous memory effects (see \citealp{distributed2024fractional}).
    \item Variable-order fractional derivatives model time-varying memory properties (see \citealp{variable2024fractional}).
    \item Hybrid fractional-delay models separate long-term memory from immediate feedback (see \citealp{hybrid2025fractional}).
    \item Distributed delay models can approximate fractional dynamics under commensurate order constraints (see \citealp{distributed2025delay}).
    \item Fractional-order experiment design under model-accuracy constraints provides minimum-power LMI inputs (Jakowluk \& Świercz, 2025), adapted here for ecological safety constraints.
\end{itemize}

\subsection{Data Analysis Protocol}
\begin{enumerate}
    \item \textbf{Preprocessing}: Detrend and normalize time series; remove outliers via median absolute deviation.
    \item \textbf{Model Fitting}: Estimate parameters for fractional, delayed, and latent models using:
          \begin{itemize}
              \item Fractional order $\alpha$ and parameters via weighted least squares in time domain.
              \item Delay $\tau$ and gain via nonlinear least squares.
              \item Latent states via expectation-maximization for linear compartmental models.
          \end{itemize}
          Model selection employs:
          \begin{itemize}
              \item Akaike Information Criterion (AIC).
              \item Bayesian Information Criterion (BIC).
              \item $k$-fold cross-validation error.
          \end{itemize}
    \item \textbf{Residual Analysis}: Inspect residual autocorrelation and partial autocorrelation for remaining structure.
    \item \textbf{Uncertainty Quantification}: Profile likelihood or Markov Chain Monte Carlo (MCMC) sampling for parameter confidence intervals.
\end{enumerate}

\subsection{Identified Gaps and Future Work}
Despite the comprehensive design, several gaps warrant further investigation:
\begin{enumerate}
    \item \textbf{Non-Stationary Memory}: The current design assumes time-invariant memory properties. Future work should incorporate time-varying fractional orders via sliding-window identification.
    \item \textbf{Spatial Heterogeneity}: The model assumes well-mixed populations; spatial extensions require reaction-diffusion fractional frameworks.
    \item \textbf{Experimental Validation of Hybrid Models}: While hybrid fractional-delay models are theoretically promising, dedicated experiments to isolate delay versus fractional components are needed.
    \item \textbf{Scaling to Multispecies Communities}: Extending the two-species predator-prey framework to multi-trophic layers remains an open challenge.
\end{itemize}
Addressing these gaps will require complementary simulation studies (e.g., bifurcation analysis of fractional-order ecological networks) and targeted follow-up experiments.

\endinput